\rhead{MA101: Calculus}
\phantomsection 
\section*{Calculus (MA101)}
\label{course:MA101}

\noindent \small{\hyperref[main:scheme]{$\leftarrow$ Back to Scheme}} \vspace{0.5cm}

\noindent \textbf{Credits:} 3 (3-0-0)

\subsection*{Course Content}
\noindent\textbf{Unit 1} \\
\textit{Differential Calculus:} Review of Single Variable Calculus: Limits, Continuity, Differentiability, Mean Value Theorems (Rolle’s, Lagrange’s, and Cauchy’s). Successive Differentiation: Leibniz’s Theorem for the nth Derivative of a Product of Functions. Series Expansions: Taylor’s and Maclaurin’s Series for Functions of a Single Variable. Applications: Curve Tracing in Cartesian and Polar Coordinates, Indeterminate Forms (L’Hôpital’s Rule).

\vspace{0.3cm}

\noindent\textbf{Unit 2} \\
\textit{Integral Calculus:} Review of Single Variable Integration: Definite Integrals as a Limit of a Sum, Fundamental Theorem of Calculus. Advanced Integration: Reduction Formulae for Trigonometric Integrals. Improper Integrals: Beta and Gamma Functions, Their Properties and Relationship. Applications of Integration: Calculating Arc Lengths, Surface Areas, and Volumes of Solids of Revolution.

\vspace{0.3cm}

\noindent\textbf{Unit 3} \\
\textit{Multivariable Differential Calculus:} Functions of Several Variables: Introduction to Functions of Two and Three Variables, Limits and Continuity. Partial Differentiation: First and Higher-Order Partial Derivatives, Total Derivative. Advanced Topics: Euler’s Theorem on Homogeneous Functions, Chain Rule for Multiple Variables. Applications: Taylor’s Series for Functions of Two Variables, Finding Maxima, Minima, and Saddle Points; Method of Lagrange Multipliers for Constrained Optimization.

\vspace{0.3cm}

\noindent\textbf{Unit 4} \\
\textit{Multivariable Integration Calculus:} Multiple Integrals: Double Integrals, Evaluation, and Change of Order of Integration. Change of Variables: Transforming Integrals from Cartesian to Polar Coordinates. Triple Integrals: Evaluation in Cartesian, Cylindrical, and Spherical Coordinates. Applications: Using Double Integrals to Compute Areas and Triple Integrals to Compute the Volume of 3D Solids.

\vspace{0.3cm}

\noindent\textbf{Unit 5} \\
\textit{3D Coordinate Geometry (Solid Geometry):} The Line: Direction Cosines and Direction Ratios, Cartesian and Vector Equations of a Line, Conditions for Coplanar Lines, Shortest Distance Between Two Skew Lines. The Plane: Cartesian and Vector Equations of a Plane, Angle Between Two Planes, Distance of a Point from a Plane. The Sphere: Equation of a Sphere, Finding Its Center and Radius, Equation of the Tangent Plane at a Point.

\newpage
\rhead{Curriculum Schema}